%% LyX 2.2.1 created this file.  For more info, see http://www.lyx.org/.
%% Do not edit unless you really know what you are doing.
\documentclass[12pt,brazil,a4paper]{article}
%\usepackage{sbc-template}
\usepackage[T1]{fontenc}
\usepackage[utf8]{inputenc}
\usepackage{geometry}
\geometry{verbose,tmargin=2cm,bmargin=2cm,lmargin=2cm,rmargin=2cm}
%\usepackage[linesnumbered,ruled,vlined,portuguese]{algorithm2e}
\usepackage{listings}
\usepackage{xcolor}
\usepackage{amsmath}

\lstset{
  basicstyle=\ttfamily\small,
  frame=single,
  breaklines=true,
  tabsize=2,
  showstringspaces=false,
  numbers=left,
  numberstyle=\tiny\color{gray},
  stepnumber=1,
  keywordstyle=\color{blue}\bfseries,
  commentstyle=\color{gray}\itshape,
  stringstyle=\color{orange!80!black}
}

% Permite marcar nomes de funcoes e tipos como palavras-chave destacadas
% no lstlisting. Uso: \addfunctions{f1,f2} e \addtypes{t1,t2} antes do bloco.
\newcommand{\addfunctions}[1]{%
  \lstset{emph={#1}, emphstyle=\bfseries}%
}
\newcommand{\addtypes}[1]{%
  \lstset{emph={[2]#1}, emphstyle={[2]\color{teal}\bfseries}}%
}

\RequirePackage{times}
\makeatletter

\newenvironment{code}
{\par\begin{list}{}{
\setlength{\rightmargin}{\leftmargin}
\setlength{\listparindent}{0pt}% needed for AMS classes
\raggedright
\setlength{\itemsep}{0pt}
\setlength{\parsep}{0pt}
\normalfont\ttfamily}%
 \item[]}
{\end{list}}

\setlength{\jot}{0.25\jot}
\usepackage{latexsym}

\makeatother

\usepackage{babel}
\begin{document}

\title{Escrevendo seu relatório}

\author{João B. Oliveira\thanks{\texttt{oliveira@inf.pucrs.br}}, Marcelo Cohen\thanks{\texttt{marcelo.cohen@pucrs.br}}\\
Escola Politécnica — PUCRS}

\date{\today}
\maketitle
\begin{abstract}
\emph{O abstract ou resumo deve dar uma descrição curta (um ou dois
parágrafos) do que será tratado no texto, para que o leitor possa
decidir se o assunto interessa ou não. Note também que a data está
em português, com o mês em minúscula. Isso é a língua portuguesa e
pode ser que seu editor de texto não saiba disso.}

Este artigo descreve alternativas de solução para o primeiro problema
proposto na disciplina de \ldots no semestre \ldots, que trata da
representação dos inteiros de 2 até 100 na forma $2^{i}\div3^{j}$.
São apresentadas três possibilidades de solução para o problema e
sua eficiência é analisada. Em seguida as representações obtidas para
os números são mostradas.

\emph{Note que seu trabalho pode pedir mais que isso! Você pode ter
que oferecer exemplos de uso, casos de teste, e outros itens que não
foram pedidos neste trabalho!}
\end{abstract}

\section*{Introdução}

\emph{O papel da introdução é descrever o problema que vai ser resolvido,
de onde ele veio, qual o seu contexto e se já sabemos algo a respeito
dele (se alguém já tentou resolver, etc.) e depois descrever os passos
que serão apresentados para a resolução, mostrando uma espécie de
roteiro do que vem a seguir. Não esqueça de que o leitor deve entender
o problema pra entender sua solução!}

Dentro do escopo da disciplina de \ldots, o primeiro problema de
\ldots pode ser resumido assim: dado um número $a=1$, outros números
podem ser produzidos seguindo-se as duas regras:
\begin{enumerate}
\item Multiplicando por 2 um número já produzido.
\item Dividindo por 3 um número já produzido e descartando a parte fracionária
(divisão inteira).
\end{enumerate}
Por exemplo, seguindo estas regras temos que
\[
10=1\times2\times2\times2\times2\times2\div3=32\div3.
\]

O problema a ser resolvido é determinar a representação de todos os
inteiros de 2 até 100 seguindo apenas estas regras. Para a saída,
os números devem ser apresentados em ordem crescente, apresentando-se
quantas vezes tivemos de usar cada regra. Segundo o enunciado do trabalho,
a linha da saída para 10 seria a seguinte:

\begin{center}
10~=~2\textasciicircum{}5~/~3\textasciicircum{}1
\par\end{center}

O enunciado também informa que os resultados possíveis não são únicos,
muitos números podem ser obtidos em mais de uma forma, portanto deve-se
dar preferência aos expoentes menores. Algumas regras e considerações
adicionais são:
\begin{enumerate}
\item Não é permitido usar inteiros especiais, como \texttt{long}, \texttt{BigInt}
e outros inteiros extra-longos. O maior tipo inteiro deve ter 32 bits.
\item É permitido usar \texttt{float,} \texttt{real} e \texttt{double} em
tarefas auxiliares, mas não é permitido calcular potências com eles.
\item Se um número não precisa ser dividido por 3, pode-se apresentar \texttt{3\textasciicircum{}0}
na saída.
\end{enumerate}
Para resolver o problema proposto, analisaremos três possíveis alternativas
de solução, bem como suas características, optando por uma soluçao
bastante eficiente ainda que pouco intuitiva. Em seguida os resultados
obtidos serão apresentados, bem como as conclusões obtidas no decorrer
do trabalho.

\emph{Talvez você tenha percebido que este texto é todo escrito no
plural: fizemos, faremos, analisaremos, apresentaremos.} \textbf{\emph{Nunca}}
\emph{usamos o singular (fiz, analisei, pensei, etc), mesmo que você
esteja trabalhando sozinho. Outra possibilidade, melhor ainda, é usar
termos passivos: faz-se, analiza-se, demonstra-se, etc.}

\emph{O texto é em espaço simples, pois é muito chato ler um trabalho
em espaço maior: fica evidente que quem escreveu precisava encher
papel. Sobre o tamanho de letra, recomenda-se 10, 11 ou 12. Este texto
foi originalmente feito em tamanho 12 porque iria ser reduzido.}

\emph{Sobre o português: não existe espaço antes de uma vírgula, depois
de dois pontos sempre vêm minúscula e coloca-se espaço em branco antes
de abrir um par de parênteses. Por dentro dos parênteses não se coloca
espaço, o texto vem colado neles. Procure estes detalhes no texto!}

\section*{Primeira solução}

Depois de considerar o problema, podemos perceber que o valor inicial
$a=1$ serve apenas para fornecer um número inicial sobre o qual podem
ser aplicadas as regras para gerar outros números, e portanto seu
valor não tem importância especial. Com esta constatação, e sabendo
que em um inteiro de 32 bits podemos armazenar números até $2^{31}$
ou até $3^{20}$ (obtivemos estes dados através de testes), podemos
propor uma solução bastante simples: bastaria criar frações $2^{i}/3^{j}$
com todas as possibilidades de expoentes para $2^{i}$ (ou seja, $0\leq i\leq31$)
e para $3^{j}$ (ou seja, $0\leq j\leq20$), verificando quais números
são gerados. O total de números a serem testados é de apenas $32*21=672$,
portanto o algoritmo deve ser executado com grande velocidade. Um
algoritmo implementando esta idéia seria parecido com este:

% Exemplo usando algorithm2e:
% https://www.sharelatex.com/learn/algorithms
%\begin{algorithm}[H]
%\SetAlgoLined
% \Para{$i \gets 1$ até $100$}{$tab[i] \gets FALSE$}
% \Para{$i \gets 1$ até $100$}{
%   \Para{$j \gets 0$ até $20$} {
%     $num \gets \frac{2^{i}}{3^{j}}$\;
%     \Se{$num <= 100$ and $tab[num] = FALSE$} {
%     $tab[num].achei =TRUE$\;
%     $tab[num].exp2 =i$\;
%     $tab[num].exp3 =j$\;
%     }
%  }
% }
% \caption{Como escrever algoritmos}
% \label{algo:exemplo}
%\end{algorithm}

\addfunctions{Algoritmo1}
\addtypes{registro,int,real}

\begin{lstlisting}
registro tab[100] {
  int achei, int exp2, int exp3
}

procedimento Algoritmo1
  para i := 1 até 100 faça
    tab[i] := FALSE
  fim
  para i := 1 até 100 faça
    para j := 0 até 20 faça
      num := ${2^i}/{3^j}$
      se num >= 100 e tab[num] = FALSE então
        tab[num].achei := TRUE
        tab[num].exp2 := i
        tab[num].exp3 := j
      fim
    fim
  fim
fim
\end{lstlisting}

Foi criada uma tabela para manter os resultados encontrados, e a cada
resultado encontrado seus expoentes são armazenados. Ao final basta
imprimir a tabela, colocando-a no formato adequado de saída (este
algoritmo é simples e não será mostrado aqui).

Também é interessante perceber que a solução apresentada até o momento
acha os menores expoentes possíveis, mas precisa manter o controle
de que números já foram encontrados, para não escrever expoentes maiores
sobre valores já armazenados. Uma maneira muito simples de descartar
este tipo de controle é fazendo o loop externo \textbf{decrescente},
garantindo assim que os menores expoentes serão encontrados por último,
e permitindo a escrita sobre valores já encontrados. Esta solução
seria:

\addfunctions{Algoritmo2}

\begin{lstlisting}
registro tab[100] {
  int exp2, int exp3
}

procedimento Algoritmo2
  para i := 31 até 0 faça
    para j := 0 até 20 faça
      num := ${2^i}/{3^j}$
      se num <= 100 então
        tab[num].exp2 := i
        tab[num].exp3 := j
      fim
    fim
  fim
fim
\end{lstlisting}

Esta idéia é bastante simples e eficiente, porém um teste mostra que
ela acha apenas 53 dos 99 números desejados. Com isto podemos concluir
que os outros 46 valores devem ser obtidos com expoentes maiores,
e os números envolvidos não podem ser representados em inteiros de
32 bits.

\emph{Note que esta seção dá uma idéia inicial, analisa-a de forma
rápida e diz porque ela não serve, também de forma rápida, sem perder
o tempo do leitor com papo furado. Essa precisão e detalhe são sinais
de um trabalho onde vem coisas mais interessantes mais tarde. Veja
que não foi usada uma linguagem de programação com todos os detalhezinhos!
Os algoritmos estão em uma linguagem simplificada, que permite que
se entenda rapidamente o que deve ser feito, sem excesso de detalhes
inúteis.}

\emph{Perceba também que o código não é cheio de explicações óbvias
e inúteis (esta linha incrementa o contador, esta linha calcula o
número), mas é clara o bastante para que outra pessoa entenda o que
é feito e reproduza o algoritmo em sua linguagem preferida. Uma falha:
terminamos dizendo que achamos 53 dos 99 números, mas sem dizer quais
foram nem falar do tempo necessário para encontrá-los. }\textbf{\emph{Moral}}\emph{:
quando terminar de falar de uma solução, apresente os resultados,
comente o tempo e as (possíveis) falhas!!}

\section*{Segunda solução}

Uma vez que a alternativa mais simples falha em quase metade dos casos,
optamos por desenvolver uma abordagem mais cuidadosa, experimentando
um truque bastante simples: se já tivermos os expoentes para certos
números (obtidos, por exemplo, com o método anterior), podemos criar
novos números a partir destes. Por exemplo, se sabemos que $2=2^{1}\div3^{0}$
e que $3=2^{5}\div3^{2}$, então podemos tentar produzir 6 através
da combinação das duas representações: $6=2*3=2^{1+5}\div3^{0+2}=2^{6}\div3^{2}.$

Infelizmente, conferindo o resultado descobrimos que $2^{6}\div3^{2}=64/9=7.1111\ldots$
e portanto não obtemos 6, e sim 7. Tabulando os resultados e prestando
atenção nos valores, temos

\vspace{0.3cm}

\begin{center}
\begin{tabular}{|c|c|l|l|}
\hline 
Valor truncado & Fórmula & Valor exato & Excesso\\
\hline 
2 & $2^{1}\div3^{0}$ & 2 & 0\\
\hline 
3 & $2^{5}\div3^{2}$ & 3.5555\ldots{} & 0.5555\ldots{}\\
\hline 
6 & $2^{6}\div3^{2}$ & 7.1111\ldots{} & 1.1111\ldots{}\\
\hline 
\end{tabular}
\par\end{center}

\vspace{0.3cm}

O fenômeno parece bastante simples: como a representação de 3 carrega
consigo um erro (excesso) de 0.5555\ldots{}, ao multiplicar 3 por
2 o erro também dobra, fazendo com que o valor obtido seja maior do
que 6. Como se pode ver na tabela, o erro para 6 é de 1.1111\ldots{},
exatamente o dobro de 0.5555\ldots{}.

Desta maneira, podemos perceber que um mecanismo simples de geração
de novos números a partir de outros que já são conhecidos teria, no
mínimo, os seguintes problemas:
\begin{enumerate}
\item São necessários valores iniciais que são usados como ``sementes''.
Estes valores iniciais podem ser obtidos com o método anterior, por
exemplo.
\item Os números sendo procurados devem ser decompostos para que se saiba
que outros números devem ser multiplicados para obtê-los, da mesma
forma que constatamos que 6 poderia ser produzido com 2 e 3. Esta
decomposição é relativamente trabalhosa, e envolve números primos
e vários tipos de testes.
\item É preciso conhecer a representação para os números primos, pois estes
não podem ser obtidos multiplicando-se outros valores. Por exemplo,
o valor 19 não poderá ser decomposto e sua representação tem que ser
descoberta de forma independente.
\item É preciso ter controle do erro sendo transmitido entre os números,
a fim de garantir que estamos chegando aos números corretos.
\end{enumerate}
Por todos estes motivos optamos por não desenvolver esta solução,
mas com o raciocínio feito até o momento podemos elaborar uma solução
aparentemente mais complexa, mas que se revela mais simples quando
analisada. Esta solução será apresentada a seguir.

\emph{Veja como esta segunda solução se apoia sobre a primeira, e
é descrita em um pouco mais de detalhe. Motivo: além de ser menos
óbvia, ela serve de ponte para a terceira alternativa, que será a
verdadeira resposta, e aqui precisamos deixar claro como funciona
o mecanismo de propagação do erro. Isto é fundamental para o que vem
a seguir. Note que esta solução é descartada com uma lista de motivos,
para que ninguém pergunte porque ela não foi implementada.}

\section*{Terceira solução}

Se já sabemos como o erro se propaga quando a regra da multiplicação
por 2 é usada, podemos pensar em uma alternativa: em vez de procurar
individualmente cada número entre 2 e 100, podemos iniciar uma busca
que vai aplicando as regras que temos e controlando o erro envolvido,
anotando os números encontrados no decorrer do processamento. É claro
que para isto precisamos saber exatamente como se propaga o erro quando
as regras são aplicadas, mas já conhecemos o caso da primeira regra,
e agora só falta analisar o que acontece quando se aplica a segunda.
A maneira mais simples é fazendo um teste, iniciando com 32 e dividindo
duas vezes, anotando os erros:

\vspace{0.3cm}

\begin{center}
\begin{tabular}{|c|c|l|l|}
\hline 
Valor truncado & Fórmula & Valor exato & Erro\\
\hline 
32 & $2^{5}\div3^{0}$ & 32 & 0\\
\hline 
10 & $2^{5}\div3^{1}$ & 10.6666\ldots{} & 0.6666\ldots{}\\
\hline 
3 & $2^{5}\div3^{2}$ & 3.5555\ldots{} & 0.5555\ldots{}\\
\hline 
\end{tabular}
\par\end{center}

\vspace{0.3cm}

É bastante claro que o erro não é simplesmente dividido por três a
cada aplicação da regra, pois se assim fosse não haveria erro para
10 ou 3, já que o erro inicial era zero. Assim, imaginamos que o erro
para a divisão seja calculado de forma mais complicada, e depois de
alguns testes a regra obtida é a seguinte: dado um número $n$ com
erro $e_{n,}$ ao aplicarmos a regra de divisão obteremos um novo
número $n/3$ com erro dado por $e_{n}/3+(n\bmod3)/3=(e_{n}+n\bmod3)/3$.
Embora esta regra pareça complicada, podemos fazer um teste com os
números anteriores: 

\vspace{0.3cm}

\begin{center}
\begin{tabular}{|c|l|l|}
\hline 
Valor & Valor exato & Erro\\
\hline 
32 & 32 & 0\\
\hline 
10 & 10.6666\ldots{} & $0/3+(32\bmod3)/3=0+2/3=0.6666\ldots$\\
\hline 
3 & 3.5555\ldots{} & $0.6666\ldots/3+(10\bmod3)/3=0.2222\ldots+0.3333\ldots=0.5555\ldots$\\
\hline 
\end{tabular}
\par\end{center}

\vspace{0.3cm}

Agora que sabemos como o erro se propaga quando as regras são aplicadas,
podemos facilmente propor uma solução que recebe um valor inicial
e aplica repetidamente as regras dadas, mantendo controle do erro
e anotando novos valores encontrados.

\emph{Perceba que a fórmula apresentada não foi provada nem justificada
de forma sólida! Isto deve ser melhorado para que um leitor entenda
}\textbf{\emph{por que}}\emph{ ela funciona! }

Este algoritmo pode funcionar de forma recursiva, da seguinte forma:

\addfunctions{pesquisa}

\begin{lstlisting}
registro tab[100] {
  int achei, int exp2, int exp3
}

para i:= 1 até 100 faça
  tab[i].achei := FALSE
fim
achados := 0

procedimento pesquisa(int num, int exp2, int exp3, real err)
  // Verificamos se o erro ultrapassa 1.
  // Se ultrapassa estamos falando de outro número...
  se err >= 1 então
    num := num + 1
    err := err - 1
  fim
  
  // Armazenamos os expoentes...
  se num <= 100 e tab[num].achei := FALSE então
    tab[num].exp2 := exp2
    tab[num].exp3 := exp3
    tab[num].achei := TRUE
    achados := achados + 1
  fim
  
  // Se ainda faltam números, reaplicar regras...
  se achados < 100 então
    pesquisa(num * 2, exp2 + 1, exp3, err * 2)
    pesquisa(num / 3, exp2, exp3 + 1, (err + n mod 3) / 3)
  fim
fim

// Primeira chamada
pesquisa(1, 0, 0, 0)
\end{lstlisting}

\emph{A solução é quase essa, mas afinal se você correr pra implementar
e funcionar de cara seria muito fácil, e isso não dá pra deixar. Num
artigo de verdade eu colocaria a solução certinha, mas este é uma
demonstração, por isso você pode pensar a respeito.}

\emph{Se você está usando linguagens orientadas a objetos: dar uma
lista de classes que você criou} \textbf{\emph{não}} \emph{é explicação
pra nada! Falar das classes não basta, você deve esclarecer o que
elas fazem! Ou seja, como trabalham, que algoritmos usam, como são
interligadas. Pense como é pouco informativo: ``\ldots{}daí fiz umas
classes A, B e C, que resolvem o problema.'' Isso diz alguma coisa?}

A primeira chamada para a função que pesquisa os números deve ser
feita com um número inicial, e é bastante simples determinar este
número: deve ser $a$, o número fornecido inicialmente para o problema.
Desta forma, passamos inicialmente o valor 1, com expoente 0 para
2 e 3 e com erro também zero. A partir daí o algoritmo é executado
até que 100 valores estejam na tabela, e em seguida a tabela pode
ser impressa no formato adequado.

Nesta versão de algoritmo não garantimos que estamos encontrando os
menores expoentes possíveis, mas o algoritmo pode ser modificado para
usar comparações com valores já existentes, ou manter listas de expoentes
em vez de apenas uma tabela.

\emph{Perceba que você não precisa ser a mãe do leitor, explicando
cada linha do código e detalhe por detalhe. Se você acha que ja explicou
o bastante e a idéia está clara, deixe que o leitor pense um pouco
para entender como as coisas funcionam. Cuidado: isso} \textbf{\emph{não}}
\emph{quer dizer que você deve deixá-lo perdido e no escuro. Você
também não pode colocar cinco páginas de código e esperar que ele
simplesmente leia e entenda.}

\section*{Resultados}

Depois de implementar o algoritmo acima em linguagem C e executá-lo
num tempo de aproximadamente meio segundo, obtivemos os seguintes
resultados:

\vspace{0.3cm}

\begin{center}
\begin{tabular}{l|l|l|l}
\texttt{1 = 2\textasciicircum{}0 / 3\textasciicircum{}0} & \texttt{26 = 2\textasciicircum{}19 / 3\textasciicircum{}9} & \texttt{51 = 2\textasciicircum{}58 / 3\textasciicircum{}33} & \texttt{76 = 2\textasciicircum{}76 / 3\textasciicircum{}44}\\
\texttt{2 = 2\textasciicircum{}1 / 3\textasciicircum{}0} & \texttt{27 = 2\textasciicircum{}46 / 3\textasciicircum{}26} & \texttt{52 = 2\textasciicircum{}39 / 3\textasciicircum{}21} & \texttt{77 = 2\textasciicircum{}57 / 3\textasciicircum{}32}\\
\texttt{3 = 2\textasciicircum{}5 / 3\textasciicircum{}2} & \texttt{28 = 2\textasciicircum{}8 / 3\textasciicircum{}2} & \texttt{53 = 2\textasciicircum{}20 / 3\textasciicircum{}9} & \texttt{78 = 2\textasciicircum{}38 / 3\textasciicircum{}20}\\
\texttt{4 = 2\textasciicircum{}2 / 3\textasciicircum{}0} & \texttt{29 = 2\textasciicircum{}16 / 3\textasciicircum{}7} & \texttt{54 = 2\textasciicircum{}66 / 3\textasciicircum{}38} & \texttt{79 = 2\textasciicircum{}19 / 3\textasciicircum{}8}\\
\texttt{5 = 2\textasciicircum{}4 / 3\textasciicircum{}1} & \texttt{30 = 2\textasciicircum{}62 / 3\textasciicircum{}36} & \texttt{55 = 2\textasciicircum{}47 / 3\textasciicircum{}26} & \texttt{80 = 2\textasciicircum{}84 / 3\textasciicircum{}49}\\
\texttt{6 = 2\textasciicircum{}9 / 3\textasciicircum{}4} & \texttt{31 = 2\textasciicircum{}24 / 3\textasciicircum{}12} & \texttt{56 = 2\textasciicircum{}9 / 3\textasciicircum{}2} & \texttt{81 = 2\textasciicircum{}65 / 3\textasciicircum{}37}\\
\texttt{7 = 2\textasciicircum{}6 / 3\textasciicircum{}2} & \texttt{32 = 2\textasciicircum{}5 / 3\textasciicircum{}0} & \texttt{57 = 2\textasciicircum{}74 / 3\textasciicircum{}43} & \texttt{82 = 2\textasciicircum{}130 / 3\textasciicircum{}78}\\
\texttt{8 = 2\textasciicircum{}3 / 3\textasciicircum{}0} & \texttt{33 = 2\textasciicircum{}13 / 3\textasciicircum{}5} & \texttt{58 = 2\textasciicircum{}55 / 3\textasciicircum{}31} & \texttt{83 = 2\textasciicircum{}46 / 3\textasciicircum{}25}\\
\texttt{9 = 2\textasciicircum{}8 / 3\textasciicircum{}3} & \texttt{34 = 2\textasciicircum{}59 / 3\textasciicircum{}34} & \texttt{59 = 2\textasciicircum{}17 / 3\textasciicircum{}7} & \texttt{84 = 2\textasciicircum{}27 / 3\textasciicircum{}13}\\
\texttt{10 = 2\textasciicircum{}5 / 3\textasciicircum{}1} & \texttt{35 = 2\textasciicircum{}21 / 3\textasciicircum{}10} & \texttt{60 = 2\textasciicircum{}82 / 3\textasciicircum{}48} & \texttt{85 = 2\textasciicircum{}92 / 3\textasciicircum{}54}\\
\texttt{11 = 2\textasciicircum{}13 / 3\textasciicircum{}6} & \texttt{36 = 2\textasciicircum{}48 / 3\textasciicircum{}27} & \texttt{61 = 2\textasciicircum{}63 / 3\textasciicircum{}36} & \texttt{86 = 2\textasciicircum{}73 / 3\textasciicircum{}42}\\
\texttt{12 = 2\textasciicircum{}10 / 3\textasciicircum{}4} & \texttt{37 = 2\textasciicircum{}10 / 3\textasciicircum{}3} & \texttt{62 = 2\textasciicircum{}44 / 3\textasciicircum{}24} & \texttt{87 = 2\textasciicircum{}54 / 3\textasciicircum{}30}\\
\texttt{13 = 2\textasciicircum{}18 / 3\textasciicircum{}9} & \texttt{38 = 2\textasciicircum{}56 / 3\textasciicircum{}32} & \texttt{63 = 2\textasciicircum{}25 / 3\textasciicircum{}12} & \texttt{88 = 2\textasciicircum{}35 / 3\textasciicircum{}18}\\
\texttt{14 = 2\textasciicircum{}7 / 3\textasciicircum{}2} & \texttt{39 = 2\textasciicircum{}18 / 3\textasciicircum{}8} & \texttt{64 = 2\textasciicircum{}71 / 3\textasciicircum{}41} & \texttt{89 = 2\textasciicircum{}16 / 3\textasciicircum{}6}\\
\texttt{15 = 2\textasciicircum{}23 / 3\textasciicircum{}12} & \texttt{40 = 2\textasciicircum{}64 / 3\textasciicircum{}37} & \texttt{65 = 2\textasciicircum{}52 / 3\textasciicircum{}29} & \texttt{90 = 2\textasciicircum{}81 / 3\textasciicircum{}47}\\
\texttt{16 = 2\textasciicircum{}4 / 3\textasciicircum{}0} & \texttt{41 = 2\textasciicircum{}45 / 3\textasciicircum{}25} & \texttt{66 = 2\textasciicircum{}33 / 3\textasciicircum{}17} & \texttt{91 = 2\textasciicircum{}146 / 3\textasciicircum{}88}\\
\texttt{17 = 2\textasciicircum{}20 / 3\textasciicircum{}10} & \texttt{42 = 2\textasciicircum{}26 / 3\textasciicircum{}13} & \texttt{67 = 2\textasciicircum{}14 / 3\textasciicircum{}5} & \texttt{92 = 2\textasciicircum{}62 / 3\textasciicircum{}35}\\
\texttt{18 = 2\textasciicircum{}9 / 3\textasciicircum{}3} & \texttt{43 = 2\textasciicircum{}53 / 3\textasciicircum{}30} & \texttt{68 = 2\textasciicircum{}79 / 3\textasciicircum{}46} & \texttt{93 = 2\textasciicircum{}43 / 3\textasciicircum{}23}\\
\texttt{19 = 2\textasciicircum{}17 / 3\textasciicircum{}8} & \texttt{44 = 2\textasciicircum{}15 / 3\textasciicircum{}6} & \texttt{69 = 2\textasciicircum{}60 / 3\textasciicircum{}34} & \texttt{94 = 2\textasciicircum{}24 / 3\textasciicircum{}11}\\
\texttt{20 = 2\textasciicircum{}44 / 3\textasciicircum{}25} & \texttt{45 = 2\textasciicircum{}80 / 3\textasciicircum{}47} & \texttt{70 = 2\textasciicircum{}41 / 3\textasciicircum{}22} & \texttt{95 = 2\textasciicircum{}89 / 3\textasciicircum{}52}\\
\texttt{21 = 2\textasciicircum{}6 / 3\textasciicircum{}1} & \texttt{46 = 2\textasciicircum{}42 / 3\textasciicircum{}23} & \texttt{71 = 2\textasciicircum{}22 / 3\textasciicircum{}10} & \texttt{96 = 2\textasciicircum{}154 / 3\textasciicircum{}93}\\
\texttt{22 = 2\textasciicircum{}14 / 3\textasciicircum{}6} & \texttt{47 = 2\textasciicircum{}23 / 3\textasciicircum{}11} & \texttt{72 = 2\textasciicircum{}68 / 3\textasciicircum{}39} & \texttt{97 = 2\textasciicircum{}70 / 3\textasciicircum{}40}\\
\texttt{23 = 2\textasciicircum{}22 / 3\textasciicircum{}11} & \texttt{48 = 2\textasciicircum{}69 / 3\textasciicircum{}40} & \texttt{73 = 2\textasciicircum{}49 / 3\textasciicircum{}27} & \texttt{98 = 2\textasciicircum{}51 / 3\textasciicircum{}28}\\
\texttt{24 = 2\textasciicircum{}30 / 3\textasciicircum{}16} & \texttt{49 = 2\textasciicircum{}31 / 3\textasciicircum{}16} & \texttt{74 = 2\textasciicircum{}30 / 3\textasciicircum{}15} & \texttt{99 = 2\textasciicircum{}32 / 3\textasciicircum{}16}\\
\texttt{25 = 2\textasciicircum{}11 / 3\textasciicircum{}4} & \texttt{50 = 2\textasciicircum{}12 / 3\textasciicircum{}4} & \texttt{75 = 2\textasciicircum{}95 / 3\textasciicircum{}56} & \texttt{100 = 2\textasciicircum{}97 / 3\textasciicircum{}57}\\
\end{tabular}
\par\end{center}

\vspace{0.3cm}

Após a obtenção destes valores, estes foram testados com ferramentas
capazes de lidar com os grandes números envolvidos, e todos os resultados
foram confirmados corretos.

\emph{Por favor, não esqueça de apresentar os resultados, afinal deve
ser isso que seu leitor quer ver!}

\section*{Conclusões}

As abordagens iniciais, mesmo oferecendo apenas resultados parciais,
contribuíram bastante para o entendimento do problema e abriram caminho
para a solução definitiva. Esta se mostrou bastante simples e eficiente,
embora tenha exigido um estudo de como o erro se propaga quando as
regras são aplicadas a novos números. A solução implementada não se
preocupou em achar sempre os menores expoentes possíveis, mas mesmo
assim ela foi bastante eficiente e os expoentes encontrados não são
excessivamente grandes (foram feitas comparações com o primeiro algoritmo
apresentado, e as diferenças foram poucas). Além disso, a solução
recursiva adotada foi razoavelmente elegante e clara, não precisando
de uma implementação complicada.

Embora não tenham sido mostrados neste artigo, foram feitos testes
com outros valores iniciais ($a\not=1)$ e para números indo até bem
mais do que 100, o algoritmo também funcionou bem.

Acreditamos ter desenvolvido uma solução interessante e barata a um
problema relativamente complexo, já que em uma linguagem de programação
usual não existe suporte direto para tratar com os números envolvidos,
e conseguimos determinar resultados corretos que envolvem grandes
números.

\emph{Suas conclusões não precisam ser geniais nem mudar o mundo.
Veja que a parte inteligente do artigo já foi feita, e veio antes!
Agora é a hora de fechar e dizer o que você acha que poderia vir depois,
etc etc. E se você está pensando em dizer coisas como ``Provavelmente
existem soluções melhores'', então diga ao menos o que você pensa
que pode melhorar (por exemplo, neste artigo eu gostaria de evitar
o uso de números em ponto flutuante e preferiria trocar por algum
tipo de controle com inteiros). Se você não der opções e só disser
que dá pra fazer melhor, está se diminundo de graça. Pra que isso?}

\emph{A bibliografia abaixo não foi usada neste artigo, e serve apenas
de exemplo para a formatação. Outros exemplos estão na página da biblioteca
da PUCRS.}
\begin{thebibliography}{1}
\bibitem{1}Brassard, G.; Bratley, P: ``\textbf{Fundamentals of algorithmics}''.
Prentice Hall, Englewood Cliffs, New Jersey, 1996.

\bibitem{2}Cormen, T.~H.; Leiserson, E.~C.; Rivest, R.~L.: \textbf{``Introduction
to Algorithms''}. Mc-Graw Hill Book Co., The MIT Electrical Engineering
and Computer Science Series, Cambridge, 1990.

\bibitem{3}Graham, R.~L.; Knuth, D.~E.; Patashnik, O.: \textbf{``Concrete
mathematics''}. Addison-Wesley, Reading, 1989.

\bibitem{4}Moret, B.~M.; Shapiro, H.~D.: ``\textbf{Algorithms from
P to NP: design and efficiency}''. Addison-Wesley, Reading, 1990.

\bibitem{5}Lueker, G.~S.: ``\emph{Some techniques for solving recurrences}''.
Computing Surveys, Vol. 12, no. 4, dezembro 1980.

\bibitem{6}Rawlins, Gregory J.~E.: ``\textbf{Compared to what? An
introduction to the analysis of algorithms}'', Computer Science Press,
New York, 1992.
\end{thebibliography}

\section*{Outros detalhes}

\emph{(Coisas que não entram num artigo, mas que vou escrever aqui
porque deve ser ditas em algum lugar). }
\begin{enumerate}
\item Em primeiríssimo lugar: você está numa Universidade e portanto deve
ter ido à escola por uns onze anos pelo menos. Você teve aulas de
português por cerca de sete desses anos, várias vezes por semana e
ainda prestou prova para o vestibular. Seu editor de texto, por outro
lado, não teve aula nenhuma e provavelmente foi (meio mal) adaptado
do inglês. Mostre um pouco de respeito por si mesmo e desligue \textbf{tudo
o que puder}. Desligue todas as correções automáticas e assuma a responsabilidade
pelo que você sabe.
\item Acentos ainda existem em português, e deverão continuar por muito
tempo. Use-os.
\item Note a consistência no que está escrito: o valor inicial $a$, por
exemplo, está sempre escrito da mesma forma, em itálico, para que
o leitor não encontre cinco tipos de letra $a$ diferentes e tenha
de lembrar o que cada uma significa. Isto mostra que o autor teve
cuidado ao escrever.
\item Note que os fontes são sempre coerentes: texto em Times, matemática
em itálico, algoritmos em Typewriter. Não há fontes maiores ou menores,
nem parágrafos com fonte diferente, margem ou espaçamento maiores
ou menores do que os outros.
\item Veja também que as potências estão sempre escritas corretamente no
texto: $x^{y}$. Não se usam notações como \emph{x\textasciicircum{}y}
porque elas são feias. Esta notação só é usada quando se fala do formato
de saída, pois ele é exigido desta forma. Da mesma forma, quando são
necessários subscritos eles devem ser escritos corretamente: $A_{i}$,
$b_{ij}$ etc. Escrever \emph{Aj} ou \emph{bij} é um mau sinal.
\item Se algo está nas referências deve estar citado no texto, e se está
citado no texto tem de aparecer nas referências!! A maneira correta
de citar uma referência é: ``Segundo \cite{3} e \cite{4}, o algoritmo
precisa\ldots{}''
\item Tenha cuidado com seu texto, ele deve transmitir a impressão de que
foi escrito com todo o cuidado do mundo. Dedique tempo a ele e tente
deixá-lo melhor, pois sua reputação (e nota) dependem disso.
\item Este artigo foi totalmente escrito usando \emph{LyX}, um dos melhores
editores de textos que existe, fazendo uso de \TeX\ e \LaTeX, que
são sem dúvida os melhores formatadores de texto do mundo. Tudo isto
é software livre rodando sob LINUX. Aprenda a usá-los, vai fazer bem.
\item Para ajudar na missão ingrata de fazer um trabalho agradável para
a leitura, este \emph{checklist} pode ajudar:

\renewcommand{\labelitemi}{\(\Box\)}
\begin{itemize}
\item Artigos não tem capa.
\item Coloquei título neste trabalho? Coloquei meu nome?
\item Usei um espaçamento legal e tamanho de fonte decente ou tem só cinco
linhas por página e as letras parecem manchetes?
\item O trabalho tem uma introdução, para que se entenda do que estou falando?
Eu contei o que ia resolver?
\item O desenvolvimento é coerente, ou parece que recortei trechos de jornal
e de meus colegas e grudei tudo?
\item Dá pra entender minhas explicações ou tem que ter poderes paranormais?
\item Tem o mínimo de código possível? Não dá pra ter ainda menos? Dá pra
não ter nenhum? Se não tem nenhum, será que deveria ter um pouco mostrando
as partes mais importantes? (Geralmente deveria)
\item Se tem código, eu fico falando durante páginas e páginas entorpecentes
e sem fim pra explicar o que ele faz, linha por linha? Posso ser \textbf{claro}
sem ser \textbf{chato}?
\item Se o trabalho pedia a resposta para algum problema, será que esqueci
de dar o resultado? Tinha que dar vários exemplos e testes?
\item Usei figuras pra esclarecer coisas que eram confusas, como aquelas
listas que usei e agora eu mesmo mal consigo entender? Como vou querer
que outra pessoa entenda?
\item Minhas figuras são claras? Mostram coisas acontecendo passo a passo
ou estão jogadas no papel? As legendas são descritivas?
\item Os dados ficam melhor se forem apresentados em tabelas? Tenho gráficos
de tempo e desempenho, se preciso? Com legendas, eixos, unidades,
ou está uma bagunça feita no Excel?
\item Lembrei de colocar minhas conclusões ao final do trabalho?
\item Coloque as minhas fontes de informação na bibliografia?
\end{itemize}
\end{enumerate}

\end{document}
